%%%%%%%%%%%%%%%%%%%%%%%%%%%%%%%%%%%%%%%%%%%%%%%%%%%%%%%%%%%%%%%%%%%%%%%%
% Chapitre 2 — État de l'art
% Inclus depuis main.tex via %%%%%%%%%%%%%%%%%%%%%%%%%%%%%%%%%%%%%%%%%%%%%%%%%%%%%%%%%%%%%%%%%%%%%%%%
% Chapitre 2 — État de l'art
% Inclus depuis main.tex via %%%%%%%%%%%%%%%%%%%%%%%%%%%%%%%%%%%%%%%%%%%%%%%%%%%%%%%%%%%%%%%%%%%%%%%%
% Chapitre 2 — État de l'art
% Inclus depuis main.tex via %%%%%%%%%%%%%%%%%%%%%%%%%%%%%%%%%%%%%%%%%%%%%%%%%%%%%%%%%%%%%%%%%%%%%%%%
% Chapitre 2 — État de l'art
% Inclus depuis main.tex via \input{chapters/related_works}
%%%%%%%%%%%%%%%%%%%%%%%%%%%%%%%%%%%%%%%%%%%%%%%%%%%%%%%%%%%%%%%%%%%%%%%%

Ce chapitre synthétise les travaux qui entourent la question centrale de la thèse : comment générer des données cliniques synthétiques utiles pour le traitement automatique des langues (TAL) médical, tout en maîtrisant les risques de réidentification.
Quatre axes sont successivement examinés : les ressources cliniques disponibles, l'augmentation de données par grands modèles de langue (GML), les techniques d'alignement et d'auto-apprentissage, puis les approches préservant la vie privée.

\subsection{Ressources cliniques et benchmarks}

Le TAL clinique s'est structuré autour de jeux de données fondateurs tels que \mimic{} et MIMIC-IV \citep{Johnson2016MIMIC-IIIDatabase}, complétés par les campagnes i2b2/n2c2 \citep{uzunerPracticalApplicationsNatural2015}.
Des initiatives récentes cherchent à dépasser le paradigme centré sur l'anglais : FRASIMED projette automatiquement des annotations en français \citep{zaghir2024frasimed}, ClinText-SP publie le plus grand corpus clinique espagnol, tandis qu'E3C-3.0 couvre plus de dix langues européennes \citep{ghosh2025e3c}.
Malgré ces avancées, le paysage reste dominé par l'anglais et par des corpus issus de traductions ou de projections, ce qui soulève des biais linguistiques et limite la validité clinique en contexte réel.
Des études multi-sites montrent en outre que les modèles entraînés sur MIMIC se dégradent fortement lors du déploiement hors de leur environnement d'origine \citep{adekkanattuEvaluatingPortabilityNLP2019}, ce qui motive les approches capables de s'adapter aux spécificités institutionnelles.

\subsection{Augmentation de données et génération par GML}

L'augmentation de données a longtemps reposé sur des techniques génériques telles que la rétro-traduction ou la substitution lexicale \citep{wei2019eda}.
L'émergence des grands modèles de langue a profondément transformé ce champ, en rendant possible la génération directe d'exemples diversifiés à partir de simples instructions.
Des méthodes telles que Self-Instruct \citep{wangSelfInstructAligningLanguage2023} exploitent ces modèles pour produire, raffiner et filtrer leurs propres données.
Dans le domaine médical, des travaux récents montrent l'intérêt de distiller des paires question-réponse synthétiques pour entraîner de petits modèles d'extraction \citep{woo2025synthetic}.

Une évolution importante concerne l'intégration d'une boucle de rétroaction : le générateur est associé à un critique, textuel, scalaire ou formel \citep{madaan2023selfrefine}.
Dans le domaine clinique, SYNFAC-EDIT corrige itérativement les erreurs factuelles des résumés générés \citep{miller2024synfacedit}, Asclepius entraîne un modèle clinique uniquement sur des notes synthétiques \citep{kweon2024asclepius}, et MedSyn associe ontologie et validation multi-passes pour le codage diagnostique \citep{kumichev2024medsyn}.
Ces approches doivent toutefois composer avec des risques propres aux pipelines auto-alimentés, tels que l'exploitation du signal de récompense et la dégénérescence progressive du modèle au fil des itérations \citep{shumailov2024modelcollapse}.

\subsection{Alignement et optimisation de préférences}

L'alignement des grands modèles de langue sur des préférences humaines ou automatiques constitue le second pilier méthodologique mobilisé dans cette thèse.
Le paradigme historique, l'apprentissage par renforcement à partir de rétroactions humaines (\gls{rlhf}), combine ajustement supervisé, modèle de récompense et optimisation de politique \citep{ouyang2022training}.
Son coût et son instabilité ont motivé des alternatives : le \gls{rlaif} remplace les annotateurs humains par des modèles juges \citep{lee2023rlaif}, tandis que l'optimisation par préférence directe (\gls{dpo}) reformule l'alignement comme un problème supervisé sur paires de préférences \citep{rafailov2023dpo}.
Des variantes ultérieures, telles que \gls{simpo}, \gls{orpo} ou \gls{kto}, suppriment le modèle de référence ou fusionnent les phases d'ajustement supervisé et d'alignement.

Une seconde famille de travaux internalise la fonction de récompense : les \emph{Self-Rewarding Language Models} optimisent leurs sorties via un juge intégré au modèle \citep{yuan2024selfrewarding}, et les approches de self-training itératif comme ReST raffinent le modèle sur ses propres générations filtrées \citep{gulcehre2023rest}.
Ces méthodes soulèvent toutefois des risques de circularité, puisqu'un même modèle produit et évalue les données, et restent majoritairement hors ligne.
Des travaux récents suggèrent que la qualité des données d'ajustement supervisé peut rivaliser avec les pipelines à base de renforcement \citep{kirk2024effects}, ce qui plaide pour des stratégies hybrides combinant génération contrôlée et alignement léger.

\subsection{Approches préservant la vie privée}

La confidentialité des données textuelles repose sur une distinction fondamentale entre anonymisation, irréversible et sortant du champ du RGPD, et pseudonymisation, qui conserve une clé de réidentification \citep{gdpr2016}.
Sur le plan technique, les pipelines cliniques ont évolué des systèmes à règles vers des architectures neuronales profondes, puis vers des systèmes hybrides déployés à grande échelle comme celui de l'\gls{aphp}, qui atteint un F1 global de \num{0.99} sur douze catégories d'entités identifiantes \citep{tannier2024aphp}.
Des travaux récents explorent en outre l'usage des grands modèles de langue eux-mêmes pour l'anonymisation \citep{staabLargeLanguageModels2025}, tout en soulignant la persistance de fuites sémantiques échappant à la substitution lexicale.

La confidentialité différentielle (\gls{dp}) offre un cadre formel pour borner le risque de divulgation \citep{dwork2014algorithmic}, mais son application au texte se heurte à la haute dimensionnalité et à la sensibilité non bornée des sorties génératives.
Plusieurs lignes de travail adaptent la \gls{dp} aux grands modèles de langue, de l'ajustement par DP-SGD à la génération de données synthétiques sous garantie formelle, jusqu'à des pipelines fédérés comme KnowledgeSG \citep{wangKnowledgeSGPrivacyPreservingSynthetic2024}, qui constitue le principal point de comparaison méthodologique pour cette thèse.

L'évaluation des données synthétiques combine trois dimensions complémentaires.
L'utilité se mesure par la performance sur des tâches en aval telles que la reconnaissance d'entités nommées ou la classification CIM.
La cohérence sémantique, longtemps estimée par \gls{bleu} ou \gls{rouge}, s'appuie désormais sur des métriques distributionnelles comme \gls{fid} ou \gls{mauve} \citep{pillutla2021mauve}, complétées par des évaluations préférentielles.
Le risque de vie privée est quant à lui sondé par des attaques empiriques : inférence d'appartenance, inférence d'attributs et attaques par liaison \citep{powar2023sok_linkage}.
Ce triptyque utilité / cohérence / risque structure le protocole d'évaluation adopté dans la suite du manuscrit.

%%%%%%%%%%%%%%%%%%%%%%%%%%%%%%%%%%%%%%%%%%%%%%%%%%%%%%%%%%%%%%%%%%%%%%%%

Ce chapitre synthétise les travaux qui entourent la question centrale de la thèse : comment générer des données cliniques synthétiques utiles pour le traitement automatique des langues (TAL) médical, tout en maîtrisant les risques de réidentification.
Quatre axes sont successivement examinés : les ressources cliniques disponibles, l'augmentation de données par grands modèles de langue (GML), les techniques d'alignement et d'auto-apprentissage, puis les approches préservant la vie privée.

\subsection{Ressources cliniques et benchmarks}

Le TAL clinique s'est structuré autour de jeux de données fondateurs tels que \mimic{} et MIMIC-IV \citep{Johnson2016MIMIC-IIIDatabase}, complétés par les campagnes i2b2/n2c2 \citep{uzunerPracticalApplicationsNatural2015}.
Des initiatives récentes cherchent à dépasser le paradigme centré sur l'anglais : FRASIMED projette automatiquement des annotations en français \citep{zaghir2024frasimed}, ClinText-SP publie le plus grand corpus clinique espagnol, tandis qu'E3C-3.0 couvre plus de dix langues européennes \citep{ghosh2025e3c}.
Malgré ces avancées, le paysage reste dominé par l'anglais et par des corpus issus de traductions ou de projections, ce qui soulève des biais linguistiques et limite la validité clinique en contexte réel.
Des études multi-sites montrent en outre que les modèles entraînés sur MIMIC se dégradent fortement lors du déploiement hors de leur environnement d'origine \citep{adekkanattuEvaluatingPortabilityNLP2019}, ce qui motive les approches capables de s'adapter aux spécificités institutionnelles.

\subsection{Augmentation de données et génération par GML}

L'augmentation de données a longtemps reposé sur des techniques génériques telles que la rétro-traduction ou la substitution lexicale \citep{wei2019eda}.
L'émergence des grands modèles de langue a profondément transformé ce champ, en rendant possible la génération directe d'exemples diversifiés à partir de simples instructions.
Des méthodes telles que Self-Instruct \citep{wangSelfInstructAligningLanguage2023} exploitent ces modèles pour produire, raffiner et filtrer leurs propres données.
Dans le domaine médical, des travaux récents montrent l'intérêt de distiller des paires question-réponse synthétiques pour entraîner de petits modèles d'extraction \citep{woo2025synthetic}.

Une évolution importante concerne l'intégration d'une boucle de rétroaction : le générateur est associé à un critique, textuel, scalaire ou formel \citep{madaan2023selfrefine}.
Dans le domaine clinique, SYNFAC-EDIT corrige itérativement les erreurs factuelles des résumés générés \citep{miller2024synfacedit}, Asclepius entraîne un modèle clinique uniquement sur des notes synthétiques \citep{kweon2024asclepius}, et MedSyn associe ontologie et validation multi-passes pour le codage diagnostique \citep{kumichev2024medsyn}.
Ces approches doivent toutefois composer avec des risques propres aux pipelines auto-alimentés, tels que l'exploitation du signal de récompense et la dégénérescence progressive du modèle au fil des itérations \citep{shumailov2024modelcollapse}.

\subsection{Alignement et optimisation de préférences}

L'alignement des grands modèles de langue sur des préférences humaines ou automatiques constitue le second pilier méthodologique mobilisé dans cette thèse.
Le paradigme historique, l'apprentissage par renforcement à partir de rétroactions humaines (\gls{rlhf}), combine ajustement supervisé, modèle de récompense et optimisation de politique \citep{ouyang2022training}.
Son coût et son instabilité ont motivé des alternatives : le \gls{rlaif} remplace les annotateurs humains par des modèles juges \citep{lee2023rlaif}, tandis que l'optimisation par préférence directe (\gls{dpo}) reformule l'alignement comme un problème supervisé sur paires de préférences \citep{rafailov2023dpo}.
Des variantes ultérieures, telles que \gls{simpo}, \gls{orpo} ou \gls{kto}, suppriment le modèle de référence ou fusionnent les phases d'ajustement supervisé et d'alignement.

Une seconde famille de travaux internalise la fonction de récompense : les \emph{Self-Rewarding Language Models} optimisent leurs sorties via un juge intégré au modèle \citep{yuan2024selfrewarding}, et les approches de self-training itératif comme ReST raffinent le modèle sur ses propres générations filtrées \citep{gulcehre2023rest}.
Ces méthodes soulèvent toutefois des risques de circularité, puisqu'un même modèle produit et évalue les données, et restent majoritairement hors ligne.
Des travaux récents suggèrent que la qualité des données d'ajustement supervisé peut rivaliser avec les pipelines à base de renforcement \citep{kirk2024effects}, ce qui plaide pour des stratégies hybrides combinant génération contrôlée et alignement léger.

\subsection{Approches préservant la vie privée}

La confidentialité des données textuelles repose sur une distinction fondamentale entre anonymisation, irréversible et sortant du champ du RGPD, et pseudonymisation, qui conserve une clé de réidentification \citep{gdpr2016}.
Sur le plan technique, les pipelines cliniques ont évolué des systèmes à règles vers des architectures neuronales profondes, puis vers des systèmes hybrides déployés à grande échelle comme celui de l'\gls{aphp}, qui atteint un F1 global de \num{0.99} sur douze catégories d'entités identifiantes \citep{tannier2024aphp}.
Des travaux récents explorent en outre l'usage des grands modèles de langue eux-mêmes pour l'anonymisation \citep{staabLargeLanguageModels2025}, tout en soulignant la persistance de fuites sémantiques échappant à la substitution lexicale.

La confidentialité différentielle (\gls{dp}) offre un cadre formel pour borner le risque de divulgation \citep{dwork2014algorithmic}, mais son application au texte se heurte à la haute dimensionnalité et à la sensibilité non bornée des sorties génératives.
Plusieurs lignes de travail adaptent la \gls{dp} aux grands modèles de langue, de l'ajustement par DP-SGD à la génération de données synthétiques sous garantie formelle, jusqu'à des pipelines fédérés comme KnowledgeSG \citep{wangKnowledgeSGPrivacyPreservingSynthetic2024}, qui constitue le principal point de comparaison méthodologique pour cette thèse.

L'évaluation des données synthétiques combine trois dimensions complémentaires.
L'utilité se mesure par la performance sur des tâches en aval telles que la reconnaissance d'entités nommées ou la classification CIM.
La cohérence sémantique, longtemps estimée par \gls{bleu} ou \gls{rouge}, s'appuie désormais sur des métriques distributionnelles comme \gls{fid} ou \gls{mauve} \citep{pillutla2021mauve}, complétées par des évaluations préférentielles.
Le risque de vie privée est quant à lui sondé par des attaques empiriques : inférence d'appartenance, inférence d'attributs et attaques par liaison \citep{powar2023sok_linkage}.
Ce triptyque utilité / cohérence / risque structure le protocole d'évaluation adopté dans la suite du manuscrit.

%%%%%%%%%%%%%%%%%%%%%%%%%%%%%%%%%%%%%%%%%%%%%%%%%%%%%%%%%%%%%%%%%%%%%%%%

Ce chapitre synthétise les travaux qui entourent la question centrale de la thèse : comment générer des données cliniques synthétiques utiles pour le traitement automatique des langues (TAL) médical, tout en maîtrisant les risques de réidentification.
Quatre axes sont successivement examinés : les ressources cliniques disponibles, l'augmentation de données par grands modèles de langue (GML), les techniques d'alignement et d'auto-apprentissage, puis les approches préservant la vie privée.

\subsection{Ressources cliniques et benchmarks}

Le TAL clinique s'est structuré autour de jeux de données fondateurs tels que \mimic{} et MIMIC-IV \citep{Johnson2016MIMIC-IIIDatabase}, complétés par les campagnes i2b2/n2c2 \citep{uzunerPracticalApplicationsNatural2015}.
Des initiatives récentes cherchent à dépasser le paradigme centré sur l'anglais : FRASIMED projette automatiquement des annotations en français \citep{zaghir2024frasimed}, ClinText-SP publie le plus grand corpus clinique espagnol, tandis qu'E3C-3.0 couvre plus de dix langues européennes \citep{ghosh2025e3c}.
Malgré ces avancées, le paysage reste dominé par l'anglais et par des corpus issus de traductions ou de projections, ce qui soulève des biais linguistiques et limite la validité clinique en contexte réel.
Des études multi-sites montrent en outre que les modèles entraînés sur MIMIC se dégradent fortement lors du déploiement hors de leur environnement d'origine \citep{adekkanattuEvaluatingPortabilityNLP2019}, ce qui motive les approches capables de s'adapter aux spécificités institutionnelles.

\subsection{Augmentation de données et génération par GML}

L'augmentation de données a longtemps reposé sur des techniques génériques telles que la rétro-traduction ou la substitution lexicale \citep{wei2019eda}.
L'émergence des grands modèles de langue a profondément transformé ce champ, en rendant possible la génération directe d'exemples diversifiés à partir de simples instructions.
Des méthodes telles que Self-Instruct \citep{wangSelfInstructAligningLanguage2023} exploitent ces modèles pour produire, raffiner et filtrer leurs propres données.
Dans le domaine médical, des travaux récents montrent l'intérêt de distiller des paires question-réponse synthétiques pour entraîner de petits modèles d'extraction \citep{woo2025synthetic}.

Une évolution importante concerne l'intégration d'une boucle de rétroaction : le générateur est associé à un critique, textuel, scalaire ou formel \citep{madaan2023selfrefine}.
Dans le domaine clinique, SYNFAC-EDIT corrige itérativement les erreurs factuelles des résumés générés \citep{miller2024synfacedit}, Asclepius entraîne un modèle clinique uniquement sur des notes synthétiques \citep{kweon2024asclepius}, et MedSyn associe ontologie et validation multi-passes pour le codage diagnostique \citep{kumichev2024medsyn}.
Ces approches doivent toutefois composer avec des risques propres aux pipelines auto-alimentés, tels que l'exploitation du signal de récompense et la dégénérescence progressive du modèle au fil des itérations \citep{shumailov2024modelcollapse}.

\subsection{Alignement et optimisation de préférences}

L'alignement des grands modèles de langue sur des préférences humaines ou automatiques constitue le second pilier méthodologique mobilisé dans cette thèse.
Le paradigme historique, l'apprentissage par renforcement à partir de rétroactions humaines (\gls{rlhf}), combine ajustement supervisé, modèle de récompense et optimisation de politique \citep{ouyang2022training}.
Son coût et son instabilité ont motivé des alternatives : le \gls{rlaif} remplace les annotateurs humains par des modèles juges \citep{lee2023rlaif}, tandis que l'optimisation par préférence directe (\gls{dpo}) reformule l'alignement comme un problème supervisé sur paires de préférences \citep{rafailov2023dpo}.
Des variantes ultérieures, telles que \gls{simpo}, \gls{orpo} ou \gls{kto}, suppriment le modèle de référence ou fusionnent les phases d'ajustement supervisé et d'alignement.

Une seconde famille de travaux internalise la fonction de récompense : les \emph{Self-Rewarding Language Models} optimisent leurs sorties via un juge intégré au modèle \citep{yuan2024selfrewarding}, et les approches de self-training itératif comme ReST raffinent le modèle sur ses propres générations filtrées \citep{gulcehre2023rest}.
Ces méthodes soulèvent toutefois des risques de circularité, puisqu'un même modèle produit et évalue les données, et restent majoritairement hors ligne.
Des travaux récents suggèrent que la qualité des données d'ajustement supervisé peut rivaliser avec les pipelines à base de renforcement \citep{kirk2024effects}, ce qui plaide pour des stratégies hybrides combinant génération contrôlée et alignement léger.

\subsection{Approches préservant la vie privée}

La confidentialité des données textuelles repose sur une distinction fondamentale entre anonymisation, irréversible et sortant du champ du RGPD, et pseudonymisation, qui conserve une clé de réidentification \citep{gdpr2016}.
Sur le plan technique, les pipelines cliniques ont évolué des systèmes à règles vers des architectures neuronales profondes, puis vers des systèmes hybrides déployés à grande échelle comme celui de l'\gls{aphp}, qui atteint un F1 global de \num{0.99} sur douze catégories d'entités identifiantes \citep{tannier2024aphp}.
Des travaux récents explorent en outre l'usage des grands modèles de langue eux-mêmes pour l'anonymisation \citep{staabLargeLanguageModels2025}, tout en soulignant la persistance de fuites sémantiques échappant à la substitution lexicale.

La confidentialité différentielle (\gls{dp}) offre un cadre formel pour borner le risque de divulgation \citep{dwork2014algorithmic}, mais son application au texte se heurte à la haute dimensionnalité et à la sensibilité non bornée des sorties génératives.
Plusieurs lignes de travail adaptent la \gls{dp} aux grands modèles de langue, de l'ajustement par DP-SGD à la génération de données synthétiques sous garantie formelle, jusqu'à des pipelines fédérés comme KnowledgeSG \citep{wangKnowledgeSGPrivacyPreservingSynthetic2024}, qui constitue le principal point de comparaison méthodologique pour cette thèse.

L'évaluation des données synthétiques combine trois dimensions complémentaires.
L'utilité se mesure par la performance sur des tâches en aval telles que la reconnaissance d'entités nommées ou la classification CIM.
La cohérence sémantique, longtemps estimée par \gls{bleu} ou \gls{rouge}, s'appuie désormais sur des métriques distributionnelles comme \gls{fid} ou \gls{mauve} \citep{pillutla2021mauve}, complétées par des évaluations préférentielles.
Le risque de vie privée est quant à lui sondé par des attaques empiriques : inférence d'appartenance, inférence d'attributs et attaques par liaison \citep{powar2023sok_linkage}.
Ce triptyque utilité / cohérence / risque structure le protocole d'évaluation adopté dans la suite du manuscrit.

%%%%%%%%%%%%%%%%%%%%%%%%%%%%%%%%%%%%%%%%%%%%%%%%%%%%%%%%%%%%%%%%%%%%%%%%

Ce chapitre synthétise les travaux qui entourent la question centrale de la thèse : comment générer des données cliniques synthétiques utiles pour le traitement automatique des langues (TAL) médical, tout en maîtrisant les risques de réidentification.
Quatre axes sont successivement examinés : les ressources cliniques disponibles, l'augmentation de données par grands modèles de langue (GML), les techniques d'alignement et d'auto-apprentissage, puis les approches préservant la vie privée.

\subsection{Ressources cliniques et benchmarks}

Le TAL clinique s'est structuré autour de jeux de données fondateurs tels que \mimic{} et MIMIC-IV \citep{Johnson2016MIMIC-IIIDatabase}, complétés par les campagnes i2b2/n2c2 \citep{uzunerPracticalApplicationsNatural2015}.
Des initiatives récentes cherchent à dépasser le paradigme centré sur l'anglais : FRASIMED projette automatiquement des annotations en français \citep{zaghir2024frasimed}, ClinText-SP publie le plus grand corpus clinique espagnol, tandis qu'E3C-3.0 couvre plus de dix langues européennes \citep{ghosh2025e3c}.
Malgré ces avancées, le paysage reste dominé par l'anglais et par des corpus issus de traductions ou de projections, ce qui soulève des biais linguistiques et limite la validité clinique en contexte réel.
Des études multi-sites montrent en outre que les modèles entraînés sur MIMIC se dégradent fortement lors du déploiement hors de leur environnement d'origine \citep{adekkanattuEvaluatingPortabilityNLP2019}, ce qui motive les approches capables de s'adapter aux spécificités institutionnelles.

\subsection{Augmentation de données et génération par GML}

L'augmentation de données a longtemps reposé sur des techniques génériques telles que la rétro-traduction ou la substitution lexicale \citep{wei2019eda}.
L'émergence des grands modèles de langue a profondément transformé ce champ, en rendant possible la génération directe d'exemples diversifiés à partir de simples instructions.
Des méthodes telles que Self-Instruct \citep{wangSelfInstructAligningLanguage2023} exploitent ces modèles pour produire, raffiner et filtrer leurs propres données.
Dans le domaine médical, des travaux récents montrent l'intérêt de distiller des paires question-réponse synthétiques pour entraîner de petits modèles d'extraction \citep{woo2025synthetic}.

Une évolution importante concerne l'intégration d'une boucle de rétroaction : le générateur est associé à un critique, textuel, scalaire ou formel \citep{madaan2023selfrefine}.
Dans le domaine clinique, SYNFAC-EDIT corrige itérativement les erreurs factuelles des résumés générés \citep{miller2024synfacedit}, Asclepius entraîne un modèle clinique uniquement sur des notes synthétiques \citep{kweon2024asclepius}, et MedSyn associe ontologie et validation multi-passes pour le codage diagnostique \citep{kumichev2024medsyn}.
Ces approches doivent toutefois composer avec des risques propres aux pipelines auto-alimentés, tels que l'exploitation du signal de récompense et la dégénérescence progressive du modèle au fil des itérations \citep{shumailov2024modelcollapse}.

\subsection{Alignement et optimisation de préférences}

L'alignement des grands modèles de langue sur des préférences humaines ou automatiques constitue le second pilier méthodologique mobilisé dans cette thèse.
Le paradigme historique, l'apprentissage par renforcement à partir de rétroactions humaines (\gls{rlhf}), combine ajustement supervisé, modèle de récompense et optimisation de politique \citep{ouyang2022training}.
Son coût et son instabilité ont motivé des alternatives : le \gls{rlaif} remplace les annotateurs humains par des modèles juges \citep{lee2023rlaif}, tandis que l'optimisation par préférence directe (\gls{dpo}) reformule l'alignement comme un problème supervisé sur paires de préférences \citep{rafailov2023dpo}.
Des variantes ultérieures, telles que \gls{simpo}, \gls{orpo} ou \gls{kto}, suppriment le modèle de référence ou fusionnent les phases d'ajustement supervisé et d'alignement.

Une seconde famille de travaux internalise la fonction de récompense : les \emph{Self-Rewarding Language Models} optimisent leurs sorties via un juge intégré au modèle \citep{yuan2024selfrewarding}, et les approches de self-training itératif comme ReST raffinent le modèle sur ses propres générations filtrées \citep{gulcehre2023rest}.
Ces méthodes soulèvent toutefois des risques de circularité, puisqu'un même modèle produit et évalue les données, et restent majoritairement hors ligne.
Des travaux récents suggèrent que la qualité des données d'ajustement supervisé peut rivaliser avec les pipelines à base de renforcement \citep{kirk2024effects}, ce qui plaide pour des stratégies hybrides combinant génération contrôlée et alignement léger.

\subsection{Approches préservant la vie privée}

La confidentialité des données textuelles repose sur une distinction fondamentale entre anonymisation, irréversible et sortant du champ du RGPD, et pseudonymisation, qui conserve une clé de réidentification \citep{gdpr2016}.
Sur le plan technique, les pipelines cliniques ont évolué des systèmes à règles vers des architectures neuronales profondes, puis vers des systèmes hybrides déployés à grande échelle comme celui de l'\gls{aphp}, qui atteint un F1 global de \num{0.99} sur douze catégories d'entités identifiantes \citep{tannier2024aphp}.
Des travaux récents explorent en outre l'usage des grands modèles de langue eux-mêmes pour l'anonymisation \citep{staabLargeLanguageModels2025}, tout en soulignant la persistance de fuites sémantiques échappant à la substitution lexicale.

La confidentialité différentielle (\gls{dp}) offre un cadre formel pour borner le risque de divulgation \citep{dwork2014algorithmic}, mais son application au texte se heurte à la haute dimensionnalité et à la sensibilité non bornée des sorties génératives.
Plusieurs lignes de travail adaptent la \gls{dp} aux grands modèles de langue, de l'ajustement par DP-SGD à la génération de données synthétiques sous garantie formelle, jusqu'à des pipelines fédérés comme KnowledgeSG \citep{wangKnowledgeSGPrivacyPreservingSynthetic2024}, qui constitue le principal point de comparaison méthodologique pour cette thèse.

L'évaluation des données synthétiques combine trois dimensions complémentaires.
L'utilité se mesure par la performance sur des tâches en aval telles que la reconnaissance d'entités nommées ou la classification CIM.
La cohérence sémantique, longtemps estimée par \gls{bleu} ou \gls{rouge}, s'appuie désormais sur des métriques distributionnelles comme \gls{fid} ou \gls{mauve} \citep{pillutla2021mauve}, complétées par des évaluations préférentielles.
Le risque de vie privée est quant à lui sondé par des attaques empiriques : inférence d'appartenance, inférence d'attributs et attaques par liaison \citep{powar2023sok_linkage}.
Ce triptyque utilité / cohérence / risque structure le protocole d'évaluation adopté dans la suite du manuscrit.
